Sources : OCDE, taux combinés d'imposition des dividendes, données 2025 ; pour la France au taux normal et pour le graphique de gauche, calcul sur les taux légaux de 2025.\par Lecture : au taux normal de 25 \%, 100 \euro{} de bénéfice deviennent 75 \euro{} après impôt sur les sociétés ; distribués, ils supportent 30 \% de prélèvement forfaitaire (22,5 \euro{}), et l'actionnaire garde 52,5 \euro{}. La grande entreprise paie en plus la contribution sociale de 3,3 \% de l'impôt ; l'actionnaire au sommet paie en plus la contribution sur les hauts revenus de 4 \%. La contribution exceptionnelle de 2025 ne concerne que les groupes dont le chiffre d'affaires dépasse 1 Md\euro{} (taux porté à 31,0 \%) ou 3 Md\euro{} (36,1 \%). Les taux des autres pays sont ceux de l'OCDE, au sommet du barème.
